\documentclass[11pt,a4paper]{ctexart}
\usepackage[utf8]{inputenc}
\usepackage{geometry}
\geometry{left=2.5cm,right=2.5cm,top=2.8cm,bottom=2.8cm}
\usepackage{amsmath,amssymb,amsfonts}
\usepackage{enumitem}
\usepackage{titlesec}
\usepackage{xcolor}
\usepackage{tcolorbox}
\usepackage{hyperref}

% 【修改点】引入画图和绝对位置放置宏包用于生成水印
\usepackage{graphicx}
\usepackage{eso-pic}

% --- 颜色定义 ---
\definecolor{mainblue}{RGB}{26, 82, 118}
\definecolor{secblue}{RGB}{41, 128, 185}
\definecolor{bglight}{RGB}{244, 246, 247}
\definecolor{textdark}{RGB}{44, 62, 80}

% --- 页面美化与字体设置 ---
\hypersetup{colorlinks=true, linkcolor=mainblue, urlcolor=secblue}
\setlist[itemize]{leftmargin=2em, itemsep=0.3em}
\setlist[enumerate]{leftmargin=2em, itemsep=0.3em}

% --- 标题样式自定义 ---
\titleformat{\section}{\large\bfseries\color{mainblue}}{\thesection}{1em}{}[{\color{mainblue}\titlerule[1pt]}]
\titleformat{\subsection}{\normalsize\bfseries\color{secblue}}{\thesubsection}{1em}{}
\titleformat{\subsubsection}{\small\bfseries\color{textdark}}{\thesubsubsection}{1em}{}

% --- 水印设置 (每页三个，斜45度，居中分布) ---
\newcommand{\WMText}{贺昱琪学习资料}
\newcommand{\WMScale}{2.28}
\newcommand{\WMAngle}{45}
\colorlet{WMColor}{black!8}

\AddToShipoutPictureBG{%
  \AtPageLowerLeft{%
    \put(\LenToUnit{0.66\paperwidth},\LenToUnit{0.70\paperheight}){%
      \rotatebox{\WMAngle}{\scalebox{\WMScale}{\textcolor{WMColor}{\WMText}}}%
    }%
    \put(\LenToUnit{0.33\paperwidth},\LenToUnit{0.50\paperheight}){%
      \rotatebox{\WMAngle}{\scalebox{\WMScale}{\textcolor{WMColor}{\WMText}}}%
    }%
    \put(\LenToUnit{0.66\paperwidth},\LenToUnit{0.30\paperheight}){%
      \rotatebox{\WMAngle}{\scalebox{\WMScale}{\textcolor{WMColor}{\WMText}}}%
    }%
  }%
}

% --- 文档信息 ---
\title{\textbf{\Huge 初中数学大纲}}
\date{\today}

% 全局设置页面样式为空，去掉所有页眉和页脚（包含页码）
\pagestyle{empty}

\begin{document}

\maketitle
% 强制取消 \maketitle 自动生成的首页页码
\thispagestyle{empty}

初中数学的知识体系由小学阶段的“具体算术与直观测量”，全面升维为“抽象代数与严密逻辑证明”。

% =========================================================================
\section{模块一：数与代数}
这是中考的基础，贯穿三年，较为简单。

% 【修改点】使用 enumerate 替代 itemize，并设置跟随章节号编号 (1.1, 1.2...)
\begin{enumerate}[label=\textbf{\thesection.\arabic*}]
    \item \textbf{数系的扩充}：
    \begin{itemize}
        \item \textbf{七年级}：引入负数，建立\textbf{有理数}体系。核心工具为数轴、绝对值（$|a| \ge 0$）与科学记数法。
        \item \textbf{八年级}：引入无理数（如 $\sqrt{2}, \pi$），完成向\textbf{实数}系的全面扩充。
    \end{itemize}
    
    \item \textbf{代数式}：
    \begin{itemize}
        \item \textbf{七年级}：\textbf{整式}的加减（单项式、多项式、同类项合并）。
        \item \textbf{八年级}：整式的乘除与\textbf{因式分解}（提公因式法、公式法、十字相乘法），并由此延伸出\textbf{分式}及其运算。
        \item \textbf{九年级}：\textbf{二次根式}的化简与混合运算。
    \end{itemize}
    
    \item \textbf{方程与不等式}：
    \begin{itemize}
        \item \textbf{七年级}：\textbf{一元一次方程}（一切方程的基础）。
        \item \textbf{八年级}：\textbf{二元一次方程组}（消元思想）、\textbf{一元一次不等式（组）}（数轴表示解集）、\textbf{分式方程}（增根与检验）。
        \item \textbf{九年级}：\textbf{一元二次方程}（降次思想：配方法、公式法求根公式，判别式 $\Delta = b^2 - 4ac$ 及韦达定理）。
    \end{itemize}
\end{enumerate}

% =========================================================================
\section{模块二：函数（重难点）}
函数研究“变量与变量之间的动态映射关系”，是历年中考压轴大题，依赖数形结合思想。

% 【修改点】这里自动会变为 2.1, 2.2...
\begin{enumerate}[label=\textbf{\thesection.\arabic*}]
    \item \textbf{八年级上：坐标系与一次函数}
    \begin{itemize}
        \item 建立平面直角坐标系，将代数与几何在直角坐标系中桥接。
        \item 研究\textbf{一次函数}（直线 $y = kx + b$），掌握待定系数法求解析式，分析斜率 $k$ 与截距 $b$ 对图像位置的决定性作用，并与一元一次方程、不等式建立联系。
    \end{itemize}
    
    \item \textbf{九年级上：反比例函数}
    \begin{itemize}
        \item 研究\textbf{反比例函数}（双曲线 $y = \frac{k}{x}$）。
        \item 重点考察 $k$ 的几何意义：图像上任意一点向两坐标轴作垂线，构成的矩形面积为 $|k|$，常与三角形、四边形结合考察综合题。
    \end{itemize}
    
    \item \textbf{九年级下：二次函数（中考压轴）}
    \begin{itemize}
        \item 研究\textbf{二次函数}（抛物线 $y = ax^2 + bx + c$）。
        \item 核心考点：顶点坐标、对称轴、最值问题。中考常将其与几何动点、线段最值、相似三角形融合，形成区分度极高的最后一道压轴题。
    \end{itemize}
\end{enumerate}

% =========================================================================
\section{模块三：空间与图形}
初中几何要求所有结论必须建立在“因为……所以……”的严密逻辑之上，格式要求严格。

% 【修改点】这里自动会变为 3.1, 3.2...
\begin{enumerate}[label=\textbf{\thesection.\arabic*}]
    \item \textbf{线与角的基础（七年级）}：相交线与平行线（同位角、内错角、同旁内角的性质与判定）。
    
    \item \textbf{三角形（几何）}：
    \begin{itemize}
        \item \textbf{七/八年级}：\textbf{全等三角形}（核心判定定理：SSS, SAS, ASA, AAS, HL）的证明，等腰/等边三角形的性质。
        \item \textbf{八年级}：\textbf{勾股定理}（直角三角形核心：$a^2 + b^2 = c^2$）及其逆定理。
        \item \textbf{九年级}：\textbf{相似三角形}（中考几何核心工具，研究图形的等比例放大缩小），以及\textbf{锐角三角函数}（$\sin, \cos, \tan$ 用于解直角三角形）。
    \end{itemize}
    
    \item \textbf{特殊四边形（八年级）}：
    \begin{itemize}
        \item 研究多边形内角和。
        \item 重点攻克平行四边形、矩形、菱形、正方形的性质与判定，理清它们之间的转化关系。
    \end{itemize}
    
    \item \textbf{圆（九年级）}：
    \begin{itemize}
        \item 研究垂径定理、圆心角与圆周角定理、切线的判定与性质。
        \item 结合正多边形计算扇形面积、弧长。
    \end{itemize}
    
    \item \textbf{图形的运动}：
    \begin{itemize}
        \item 轴对称、平移、旋转。复杂几何压轴题常常需要通过“旋转”或“作对称点（如将军饮马模型）”来构造辅助线，破解难题。
    \end{itemize}
\end{enumerate}

% =========================================================================
\section{模块四：统计与概率}


% 【修改点】这里自动会变为 4.1, 4.2...
\begin{enumerate}[label=\textbf{\thesection.\arabic*}]
    \item \textbf{统计学（八年级）}：
    \begin{itemize}
        \item 了解抽样调查与普查的适用场景，掌握总体、个体、样本容量的概念。
        \item 集中趋势评估：平均数、中位数、众数。
        \item 离散程度评估：\textbf{方差}（$s^2$）与标准差，用于衡量数据的波动大小和稳定性（方差越小越稳定）。
    \end{itemize}
    
    \item \textbf{概率（九年级）}：
    \begin{itemize}
        \item 区分必然事件、不可能事件与随机事件。
        \item 掌握使用\textbf{列表法}和\textbf{树状图法}穷举所有等可能结果，精确计算两步及以上随机事件发生的概率。
    \end{itemize}
\end{enumerate}

\end{document}